% ---------------------------------------------------------------------------
% Introduction (Section 1). Citation keys match report/references.bib.
% ---------------------------------------------------------------------------
\section{Introduction}
\label{sec:intro}

Spoken digits and a spoken ``yes'' or ``no'' are enough to run a mobile-money transfer, a health hotline or a
phone menu. For users with limited literacy, voice is often the only practical interface~\cite{doumbouya2021using}.
Swahili is widely spoken across East Africa. Yet it remains a low-resource language for speech technology: it has
only limited labelled data~\cite{joshi2020state}, and it was missing from the first releases of the largest open
crowdsourced speech corpus~\cite{ardila2020common}. Recognising a small Swahili vocabulary reliably, cheaply and
with honest confidence is therefore a useful problem in its own right, as well as a test of how well modern sequence
models transfer to an African language.

We study the Zindi Swahili Audio Classification challenge~\cite{zindi2021swahili}: 4{,}200 crowdsourced clips of
12 words, the digits \textit{moja} to \textit{kumi} (one to ten), \textit{ndio} (yes) and \textit{hapana} (no),
scored by log loss. A spoken word is a \emph{sequence} of acoustic frames, and its identity depends on the order of
its sounds. The vocabulary makes this concrete: \textit{sita} (six) and \textit{tisa} (nine) contain the same sounds
in a different order. Our research question is therefore:
\begin{quote}
  \emph{How effectively can sequential modelling approaches recognise spoken Swahili words, and what evidence
  supports the strengths and limitations of each approach?}
\end{quote}

We answer it with five approaches that differ in \emph{how} they model temporal order. A1 discards order: it is a
logistic regression on MFCC statistics. A2 models order with a Markov chain: a GMM-HMM per word. A3 uses
recurrence: a bidirectional LSTM with attention. A4 uses convolution over time: TC-ResNet. A5 uses self-attention
after multilingual self-supervised pretraining: a fine-tuned XLS-R. All five share one frozen data split, one feature
pipeline and one evaluation code base. They were tuned on validation only, in 55 logged development runs, and the
test split was used only after every configuration was frozen.

Our contributions are:
\begin{itemize}
  \item \emph{A controlled comparison of five ways of modelling order} on an African-language speech task. Neural
        models are trained with three seeds, and the test results carry bootstrap confidence intervals and
        Holm-corrected McNemar tests (Section~\ref{sec:results}).
  \item \emph{Direct evidence of how the models use order.} A temporal-order stress test reverses or shuffles each
        word. It shows that all order-aware models depend on order, but that the pretrained A5 relies mainly on
        order within about 100~ms (Section~\ref{sec:errors}).
  \item \emph{A data-size curve.} With only 24 training clips per word, the pretrained model is already about as
        accurate as the models trained from scratch on all 245.
  \item \emph{An error analysis without a human listener.} It combines confident learning with an external Swahili
        recogniser. The clips that defeat every order-aware model are far more often quiet, noisy or badly cropped,
        and include a likely label error and words said in English.
  \item \emph{An open, reproducible pipeline.} The notebooks rebuild every table and figure in this report from the
        logged runs and predictions.
\end{itemize}

Section~\ref{sec:related} reviews related work, Section~\ref{sec:data} the data, and Section~\ref{sec:method} the
methods and metrics. Section~\ref{sec:results} reports and discusses the results, Section~\ref{sec:errors} analyses
the errors and limitations, and Section~\ref{sec:conclusion} concludes.
